\answer{ex:13:1}{(a)~$q=100^{1/99}\approx1.048$, paso de $\approx4.8\%$; rango de $\approx2.2\cdot10^3$ ($67$~dB). (b)~Paso de $22f_1$, es decir, el $220\%$ con $10f_1$.}
\answer{ex:13:2}{$P_{\text{fallo}}\approx1.8\cdot10^{-4}$ y $2.8\cdot10^{-16}$: unos $300$ fallos por día con $S=2$ y $5\cdot10^{-10}$ por día con $S=4$.}
\answer{ex:13:3}{(a)~$H(s)=eF_1T_c/(1+sT_c)^2$, filtro de segundo orden con frecuencia de corte $1/(2\pi T_c)\approx3.2$~Hz. (b)~$r\approx22$~espigas/s (por el primer armónico, $\approx20$).}
\answer{ex:13:4}{(a)~Con $Gd=1/e$ la raíz $s=-1/d$ es doble, y con ningún otro $G$ la raíz más a la derecha queda más a la izquierda. (b)~$d=30\ms$: $G\approx12$~s$^{-1}$, $\tau=30\ms$; $d=150\ms$: $G\approx2.5$~s$^{-1}$, $\tau=150\ms$; la constante de tiempo es igual al retardo.}
\answer{ex:13:5}{(a)~$E\approx0.427$, $T\approx1.70\,\tau_a=0.68$~s (el modelo con $\tau=20\ms$ da $0.84$~s). (b)~La ecuación no contiene $I$; el ritmo existe con $b>\beta-1$.}
\answer{ex:13:6}{Con $b=0.8$ no hay ritmo; con $b=1.05$, $1.2$, $1.5$, $2$, $4$, $T\approx2.73$, $1.74$, $1.19$, $0.84$, $0.45$~s; con cualquier $I$, $T=0.840$~s: la entrada solo cambia la amplitud, no el ritmo.}
\answer{ex:13:7}{La sustitución $s=u-a$ da la máxima tasa de amortiguamiento $a+1/d$ con $G=e^{-1-ad}/d$; $a\ge33.3$~s$^{-1}$, $G=e^{-2}/d\approx4.5$~s$^{-1}$; cuanto mayor la cocontracción (mayor $a$), menor debe ser $G$.}
\answer{ex:13:8}{Problema abierto. Por ejemplo, un umbral en la carga de fatiga $b[r_i-r_0]_+$ da $I_{\min}=\beta b r_0/(b-\beta+1)\approx0.53$ con $b=4$, $r_0=0.2$, y un período de $1.8$~s con $I=0.6$ a $0.52$~s con $I=3$.}
